% ============================================================
\chapter*{Pr\'eface}
\addcontentsline{toc}{chapter}{Pr\'eface}
% ============================================================

\section*{Objectifs du cours}

Ce cours de \textbf{Processus Stochastiques} s'adresse aux \'etudiants de niveau
Master en math\'ematiques, physique th\'eorique, finance quantitative et sciences
de l'ing\'enieur. Il vise \`a fournir une pr\'esentation rigoureuse, fond\'ee sur la
th\'eorie de la mesure, des principaux processus al\'eatoires et de leurs applications.

\medskip

L'objectif principal est de conduire l'\'etudiant depuis les rappels de probabilit\'es
jusqu'\`a la construction de l'int\'egrale stochastique d'It\^o et \`a la r\'esolution
d'\'equations diff\'erentielles stochastiques (EDS), en passant par l'\'etude approfondie
des cha\^ines de Markov, du processus de Poisson, des martingales et du mouvement
brownien.

\medskip

\`A l'issue de ce cours, l'\'etudiant sera capable de~:
\begin{itemize}
  \item Ma\^itriser le formalisme mesure-th\'eorique des processus stochastiques.
  \item Analyser les cha\^ines de Markov en temps discret et continu~: classifier
    les \'etats, calculer les distributions stationnaires, \'etablir l'ergodicit\'e.
  \item Manipuler les martingales~: appliquer les in\'egalit\'es fondamentales, les
    th\'eor\`emes de convergence et le th\'eor\`eme d'arr\^et optionnel.
  \item Comprendre les propri\'et\'es fines du mouvement brownien.
  \item Construire l'int\'egrale stochastique d'It\^o et appliquer la formule d'It\^o.
  \item R\'esoudre des \'equations diff\'erentielles stochastiques et les appliquer
    en finance, en physique et en biologie.
\end{itemize}

\section*{Pr\'erequis}

Les pr\'erequis essentiels sont~:
\begin{itemize}
  \item \textbf{Th\'eorie de la mesure et int\'egration}~: tribus, mesures, int\'egrale
    de Lebesgue, th\'eor\`emes de convergence (domin\'ee, monotone), espaces $L^p$.
  \item \textbf{Probabilit\'es de niveau licence}~: variables al\'eatoires, esp\'erance,
    variance, loi des grands nombres, th\'eor\`eme central limite.
  \item \textbf{Analyse fonctionnelle \'el\'ementaire}~: espaces de Banach et de Hilbert,
    convergences, compl\'etude.
  \item \textbf{Alg\`ebre lin\'eaire}~: matrices, valeurs propres, d\'ecomposition
    spectrale.
  \item \textbf{\'Equations diff\'erentielles ordinaires}~: th\'eor\`emes d'existence
    et d'unicit\'e, m\'ethodes de r\'esolution.
\end{itemize}

\section*{Organisation du cours}

Le cours est organis\'e en onze chapitres, suivant une progression logique~:

\begin{enumerate}
  \item \textbf{Rappels de probabilit\'es}~: espace de probabilit\'e, esp\'erance
    conditionnelle, modes de convergence. Ce chapitre fournit les fondations
    mesure-th\'eoriques n\'ecessaires \`a la suite du cours.

  \item \textbf{Processus stochastiques~--- D\'efinitions}~: d\'efinition formelle
    d'un processus stochastique, filtrations, temps d'arr\^et, processus
    adapt\'es et pr\'evisibles. Ces notions constituent le langage de base.

  \item \textbf{Cha\^ines de Markov en temps discret}~: propri\'et\'e de Markov,
    classification des \'etats, th\'eor\`emes ergodiques, mesures stationnaires.

  \item \textbf{Cha\^ines de Markov en temps continu}~: g\'en\'erateur infinit\'esimal,
    \'equations de Kolmogorov, processus de naissance et de mort.

  \item \textbf{Processus de Poisson}~: construction axiomatique, propri\'et\'es
    fondamentales, processus de Poisson compos\'e et non homog\`ene.

  \item \textbf{Martingales en temps discret}~: d\'efinition, in\'egalit\'es
    fondamentales (Doob), th\'eor\`emes de convergence, th\'eor\`eme d'arr\^et optionnel.

  \item \textbf{Martingales en temps continu}~: r\'egularit\'e des trajectoires,
    d\'ecomposition de Doob--Meyer, martingales locales.

  \item \textbf{Mouvement brownien}~: construction (L\'evy--Ciesielski),
    propri\'et\'es des trajectoires, non-d\'erivabilit\'e, loi du logarithme it\'er\'e,
    principe de r\'eflexion.

  \item \textbf{Int\'egrale stochastique d'It\^o}~: construction par processus
    \'etagn\'es, extension isom\'etrique aux processus $L^2$, propri\'et\'es de martingale.

  \item \textbf{Formule d'It\^o et EDS}~: formule de changement de variables d'It\^o,
    formule d'It\^o multidimensionnelle, existence et unicit\'e des solutions fortes.

  \item \textbf{Applications}~: finance (mod\`ele de Black--Scholes, couverture),
    physique (\'equation de Langevin, diffusion), biologie (mod\`eles de population
    stochastiques).
\end{enumerate}

\section*{Conventions et notations}

Tout au long de ce cours, nous adoptons les conventions suivantes~:

\begin{itemize}
  \item $(\Omega, \mathcal{F}, \PP)$ d\'esigne un espace de probabilit\'e fix\'e.
  \item $\{\mathcal{F}_t\}_{t \geq 0}$ ou $\{\mathcal{F}_n\}_{n \geq 0}$ d\'esigne
    une filtration satisfaisant les conditions habituelles (compl\'etude, continuit\'e \`a droite).
  \item $\E[X]$ d\'esigne l'esp\'erance de la variable al\'eatoire $X$.
  \item $\E[X \mid \mathcal{G}]$ d\'esigne l'esp\'erance conditionnelle de $X$ sachant
    la tribu $\mathcal{G}$.
  \item $\PP(A)$ d\'esigne la probabilit\'e de l'\'ev\'enement $A$.
  \item $\Var(X)$ et $\Cov(X, Y)$ d\'esignent la variance et la covariance.
  \item $\R$, $\N$, $\Z$, $\C$ d\'esignent les ensembles des r\'eels, naturels,
    entiers et complexes. On note $\R_+ = [0, +\infty)$.
  \item $\norm{\cdot}$ et $\abs{\cdot}$ d\'esignent respectivement une norme et la
    valeur absolue.
  \item $\mathbf{1}_A$ d\'esigne la fonction indicatrice de l'ensemble $A$.
  \item p.s.\ signifie <<~presque s\^urement~>>, v.a.\ <<~variable al\'eatoire~>>,
    EDS <<~\'equation diff\'erentielle stochastique~>>.
  \item $\xrightarrow{\text{p.s.}}$, $\xrightarrow{L^p}$, $\xrightarrow{\PP}$,
    $\xrightarrow{\mathcal{L}}$ d\'esignent les convergences presque s\^ure, dans $L^p$,
    en probabilit\'e et en loi.
\end{itemize}

\section*{Approche p\'edagogique}

Chaque chapitre contient~:
\begin{itemize}
  \item Des \textbf{d\'efinitions} pr\'ecises dans le cadre de la th\'eorie de la mesure.
  \item Des \textbf{th\'eor\`emes} avec d\'emonstrations compl\`etes ou esquiss\'ees.
  \item Des \textbf{exemples} d\'etaill\'es pour illustrer les concepts abstraits.
  \item Des \textbf{exercices} de difficult\'e vari\'ee, certains avec indications.
  \item Des encadr\'es \og Intuition \fg pour d\'evelopper l'intuition probabiliste.
  \item Des encadr\'es \og Attention \fg pour signaler les pi\`eges classiques.
  \item Des encadr\'es \og Formules cl\'es \fg r\'ecapitulant les r\'esultats essentiels.
\end{itemize}

\section*{R\'ef\'erences bibliographiques}

Les ouvrages suivants ont inspir\'e ce cours~:

\begin{itemize}
  \item \textsc{Br\'emaud}, P. \emph{Markov Chains: Gibbs Fields, Monte Carlo Simulation
    and Queues}. Springer, 1999.
  \item \textsc{Durrett}, R. \emph{Probability: Theory and Examples}. Cambridge
    University Press, 5\`eme \'ed., 2019.
  \item \textsc{Karatzas}, I. et \textsc{Shreve}, S.E. \emph{Brownian Motion and
    Stochastic Calculus}. Springer, 2\`eme \'ed., 1991.
  \item \textsc{Le Gall}, J.-F. \emph{Mouvement brownien, martingales et calcul
    stochastique}. Springer, 2013.
  \item \textsc{Norris}, J.R. \emph{Markov Chains}. Cambridge University Press, 1997.
  \item \textsc{{\O}ksendal}, B. \emph{Stochastic Differential Equations}. Springer,
    6\`eme \'ed., 2003.
  \item \textsc{Revuz}, D. et \textsc{Yor}, M. \emph{Continuous Martingales and
    Brownian Motion}. Springer, 3\`eme \'ed., 1999.
  \item \textsc{Shreve}, S.E. \emph{Stochastic Calculus for Finance~I \& II}.
    Springer, 2004.
  \item \textsc{Williams}, D. \emph{Probability with Martingales}. Cambridge
    University Press, 1991.
\end{itemize}

\section*{Remerciements}

L'auteur remercie les coll\`egues et \'etudiants qui, par leurs remarques et
suggestions, ont contribu\'e \`a am\'eliorer ce cours au fil des ann\'ees. Les
discussions avec les praticiens de la finance quantitative et les physiciens
ont \'egalement enrichi la pr\'esentation des applications.

\vfill

\begin{flushright}
\textit{L'auteur}\\
Mars 2026
\end{flushright}
